% GENERATED FILE. Edit canonical lessons, then run npm run generate:books.
% canonical-source-hash: fnv1a64:c923520a76b2463c
% canonical-lessons: ZH-C365-jianyi, ZH-C365-keshi, ZH-W365-ke-can, ZH-C365-kejian, ZH-C365-buru, ZH-C365-zhonggao

\chapter{Yes, but}
\label{ch:zh-yes-but}
\hlchaptermodality{365}

\begin{chapteropening}
\textbf{By the end of this chapter:} \emph{I can say but, suggest what someone should do and say what goes to show.}

Complete the last lesson of chapter 365: I can say but, suggest what someone should do and say what goes to show.
\end{chapteropening}

\section[jiànyì]{jiànyì — to suggest}

\label{lesson:ZH-C365-jianyi}



\begin{quote}
Before the new one: say the Mandarin for to give official approval, then the Mandarin for to remind someone.
\end{quote}

\subsection*{You'll want to know: jiànyì}
\emph{jiànyì} — \textquotedblleft{}to suggest\textquotedblright{}.

\begin{quote}
\emph{yīnwèi nǐ bìng le suǒyǐ wǒ jiànyì nǐ qù kàn yīshēng} — You're ill, so I suggest you go and see a doctor.
\end{quote}

Its characters come later; for now it is a word for the ear.

\begin{cousinweb}
\emph{jiànyì} is built from \emph{jiàn} and \emph{yì}, both inherited from Middle Chinese, \emph{kjonH} and \emph{ngjeH} (Guangyun, in Baxter's transcription), which glosses the first 'to set up' and the second 'to deliberate'.
\end{cousinweb}

\begin{grammarlens}[title={What you've built}]
One more word for giving reasons.
\end{grammarlens}

\subsection*{Your turn}
\begin{itemize}
\raggedright
  \item \emph{Say it:} \emph{jiànyì}
  \item \emph{Say it:} \emph{jiànyì}, once more
  \item \emph{Say it:} \emph{yīnwèi nǐ bìng le suǒyǐ wǒ jiànyì nǐ qù kàn yīshēng} — You're ill, so I suggest you go and see a doctor.
  \item \emph{Recall:} say the Mandarin for to give official approval, then the Mandarin for to remind someone, then say \emph{jiànyì} again
\end{itemize}

\begin{tcolorbox}[breakable,colback=teal!4,colframe=teal!35!black,title={Before you move on}]
What does jiànyì mean? (\textquotedblleft{}To suggest\textquotedblright{}.) Say it once more.
\end{tcolorbox}
\section[\texorpdfstring{\zh{可是} (kěshì)}{可是 (kěshì)}]{\zh{可是} (kěshì) — but; yet}

\label{lesson:ZH-C365-keshi}



\begin{quote}
Before the new one: say the Mandarin for to remind someone, then the Mandarin for to suggest.
\end{quote}

\subsection*{You'll want to know: \zh{可是}}
\textbf{\zh{可是}} — \emph{kěshì} — \textquotedblleft{}but; yet\textquotedblright{}.

\begin{quote}
\textbf{\zh{你说得很有道理可是我不想去}} \emph{nǐ shuō de hěn yǒu dàoli kěshì wǒ bù xiǎng qù} — You've got a point, but I don't want to go.
\end{quote}

The new character, in pieces: \zh{可} — five strokes: a long bar on top, \zh{口} \emph{kǒu} under it on the left, three strokes, and a long vertical down the right that hooks at the foot. The next lesson puts it in your hand.

\begin{cousinweb}
\textbf{\zh{可是}} (\emph{kěshì}) is built from \emph{kě} and \zh{是}, both inherited from Middle Chinese, \emph{khaX} and \emph{dzyeX} (Guangyun, in Baxter's transcription), which glosses the first 'to permit' and the second 'right, as against wrong'. How the pair came to mean but is not settled.
\end{cousinweb}

\begin{grammarlens}[title={What you've built}]
One more word for giving reasons.
\end{grammarlens}

\subsection*{Your turn}
\begin{itemize}
\raggedright
  \item \emph{Say it:} \emph{kěshì}
  \item \emph{Say it:} \emph{kěshì}, once more
  \item \emph{Say it:} \emph{nǐ shuō de hěn yǒu dàoli kěshì wǒ bù xiǎng qù} — You've got a point, but I don't want to go.
  \item \emph{Recall:} say the Mandarin for to remind someone, then the Mandarin for to suggest, then say \emph{kěshì} again
\end{itemize}

\begin{tcolorbox}[breakable,colback=teal!4,colframe=teal!35!black,title={Before you move on}]
What does \zh{可是} mean? (\textquotedblleft{}But; yet\textquotedblright{}.) Say it once more.
\end{tcolorbox}
\section[\texorpdfstring{\zh{可} (kě)}{可 (kě)}]{\zh{可} — a bar, a mouth and a hook}

\label{lesson:ZH-W365-ke-can}



\begin{quote}
The word first, then the character inside it.

\begin{itemize}
\raggedright
  \item \emph{Recall:} say \emph{kěshì} — \textbf{R1}, one lesson back, before it decays
\end{itemize}

\begin{quote}
\zh{口} \zh{是}
\end{quote}

\emph{Read it:} each one aloud, with nothing said first — \emph{kǒu}, \emph{shì}
\end{quote}

\subsection*{Script — a bar, a mouth and a hook}
\hlblockfigure{\detokenize{figures/ZH-W365-ke-can-filmstrip.pdf}}{How \zh{可} is written, stroke by stroke}

\begin{quote}
\zh{可}
\end{quote}

\textbf{One piece, five strokes.} Count it in three parts.

\textbf{A long horizontal} across the top.

\textbf{\zh{口}, three strokes}, small, under the bar on the left: a vertical, a horizontal that turns down, and the bottom horizontal.

\textbf{A long vertical down the right side}, from the bar to the foot, hooking up to the left.

1 + 3 + 1 = five.

\subsection*{Your turn}
\begin{itemize}
\raggedright
  \item \emph{Write it:} \zh{可}, all five, with the model in front of you
  \item \emph{Cover:} write it again from memory, and count the pieces, not the strokes
\end{itemize}

\begin{tcolorbox}[breakable,colback=teal!4,colframe=teal!35!black,title={Before you move on}]
How many strokes in \zh{可}? (\textbf{Five}.) What is the small box? (\textbf{\zh{口}}.) What comes last? (\textbf{The long vertical with its hook}.)

Source: \href{https://raw.githubusercontent.com/chanind/hanzi-writer-data/68d10a4b21150cae5e1ebbd223eed289cf32d90c/data/\%E5\%8F\%AF.json}{Hanzi Writer Data \zh{可}, pinned stroke paths and medians}.
\end{tcolorbox}
\section[\texorpdfstring{\zh{可见} (kějiàn)}{可见 (kějiàn)}]{\zh{可见} (kějiàn) — which goes to show; so it is clear that}

\label{lesson:ZH-C365-kejian}



\begin{quote}
Before the new one: say the Mandarin for to suggest, then the Mandarin for but (in everyday speech).
\end{quote}

\subsection*{You'll want to know: \zh{可见}}
\textbf{\zh{可见}} — \emph{kějiàn} — \textquotedblleft{}which goes to show; so it is clear that\textquotedblright{}.

\begin{quote}
\textbf{\zh{他每天都准时来可见他很喜欢这儿}} \emph{tā měitiān dōu zhǔnshí lái kějiàn tā hěn xǐhuan zhèr} — He turns up on time every day, which shows he really likes it here.
\end{quote}

\begin{cousinweb}
\textbf{\zh{可见}} (\emph{kějiàn}) is built from \zh{可} 'can' and \zh{见} 'to see', both inherited from Middle Chinese, \emph{khaX} and \emph{kenH} (Guangyun, in Baxter's transcription), which glosses the first 'to permit' and the second 'to see'.
\end{cousinweb}

\begin{grammarlens}[title={What you've built}]
One more word for giving reasons.
\end{grammarlens}

\subsection*{Your turn}
\begin{itemize}
\raggedright
  \item \emph{Say it:} \emph{kějiàn}
  \item \emph{Say it:} \emph{kějiàn}, once more
  \item \emph{Say it:} \emph{tā měitiān dōu zhǔnshí lái kějiàn tā hěn xǐhuan zhèr} — He turns up on time every day, which shows he really likes it here.
  \item \emph{Recall:} say the Mandarin for to suggest, then the Mandarin for but (in everyday speech), then say \emph{kějiàn} again
\end{itemize}

\begin{tcolorbox}[breakable,colback=teal!4,colframe=teal!35!black,title={Before you move on}]
What does \zh{可见} mean? (\textquotedblleft{}Which goes to show; so it is clear that\textquotedblright{}.) Say it once more.
\end{tcolorbox}
\section[\texorpdfstring{\zh{不如} (bùrú)}{不如 (bùrú)}]{\zh{不如} (bùrú) — to do better to; might as well; to be not as good as}

\label{lesson:ZH-C365-buru}



\begin{quote}
Before the new one: say the Mandarin for but (in everyday speech), then the Mandarin for which goes to show.
\end{quote}

\subsection*{You'll want to know: \zh{不如}}
\textbf{\zh{不如}} — \emph{bùrú} — \textquotedblleft{}to do better to; might as well; to be not as good as\textquotedblright{}.

\begin{quote}
\textbf{\zh{我们不如走路去}} \emph{wǒmen bùrú zǒulù qù} — We might as well walk.
\end{quote}

The tone of bù stays falling here, because the next syllable is not a falling tone.

\begin{cousinweb}
\textbf{\zh{不如}} (\emph{bùrú}) is built from \zh{不} 'not' and \zh{如} 'like', both inherited from Middle Chinese, \emph{pjut} and \emph{nyo} (Guangyun, in Baxter's transcription), which glosses the second 'like' and 'as'.
\end{cousinweb}

\begin{grammarlens}[title={What you've built}]
One more word for giving reasons.
\end{grammarlens}

\subsection*{Your turn}
\begin{itemize}
\raggedright
  \item \emph{Say it:} \emph{bùrú}
  \item \emph{Say it:} \emph{bùrú}, once more
  \item \emph{Say it:} \emph{wǒmen bùrú zǒulù qù} — We might as well walk.
  \item \emph{Recall:} say the Mandarin for but (in everyday speech), then the Mandarin for which goes to show, then say \emph{bùrú} again
\end{itemize}

\begin{tcolorbox}[breakable,colback=teal!4,colframe=teal!35!black,title={Before you move on}]
What does \zh{不如} mean? (\textquotedblleft{}To do better to; might as well; to be not as good as\textquotedblright{}.) Say it once more.
\end{tcolorbox}
\section[zhōnggào]{zhōnggào — word of advice; warning}

\label{lesson:ZH-C365-zhonggao}



\begin{quote}
Before the new one: say the Mandarin for which goes to show, then the Mandarin for might as well.
\end{quote}

\subsection*{You'll want to know: zhōnggào}
\emph{zhōnggào} — \textquotedblleft{}word of advice; warning\textquotedblright{}.

\begin{quote}
\emph{nǐ yīnggāi tīng tā de zhōnggào} — You should take his advice.
\end{quote}

Its characters come later; for now it is a word for the ear.

\begin{cousinweb}
\emph{zhōnggào} is built from \emph{zhōng} and \emph{gào}, both inherited from Middle Chinese, \emph{trjuwng} and \emph{kawH} (Guangyun, in Baxter's transcription), which glosses the first 'selfless' and the second 'to report'.
\end{cousinweb}

\begin{grammarlens}[title={What you've built}]
One more word for giving reasons.
\end{grammarlens}

\subsection*{Your turn}
\begin{itemize}
\raggedright
  \item \emph{Say it:} \emph{zhōnggào}
  \item \emph{Say it:} \emph{zhōnggào}, once more
  \item \emph{Say it:} \emph{nǐ yīnggāi tīng tā de zhōnggào} — You should take his advice.
  \item \emph{Recall:} say the Mandarin for which goes to show, then the Mandarin for might as well, then say \emph{zhōnggào} again
\end{itemize}

\begin{tcolorbox}[breakable,colback=teal!4,colframe=teal!35!black,title={Before you move on}]
What does zhōnggào mean? (\textquotedblleft{}Word of advice; warning\textquotedblright{}.) Say it once more.
\end{tcolorbox}
